% Class-compatibility preamble retained for a reader excerpt of
%   Sahar Diskin, Anna Geisler, "Perfect Matching in Product Graphs and in their
%   Random Subgraphs", arXiv:2404.14020v4, CC BY 4.0.
%
% The source pins its own publisher class (scrartcl); the delivery wrapper accepts
% only the classes article and amsart, so this file supplies the amsart baseline and
% then copies the author's own declaration block mechanically and unchanged
% (pinned-original-source.tex, source lines 30-57: the author's newtheorem
% declarations and the author's macro declarations).  The author's own package
% list of source lines 4-21 is reused only for the entries the retained text needs
% (enumitem, amsmath, amssymb, amsthm, mathtools, xcolor, geometry) and the
% cross-reference packages already needed by the author's own \Cref/\eqref usage
% (hyperref, cleveref); tikz, caption, subcaption, float, graphicx, wrapfig and
% abstract serve figures and the title block that this excerpt does not carry.
% No mathematical macro, statement, symbol or assumption is introduced, changed
% or removed by this file.
\documentclass[10pt,a4paper]{amsart}
\usepackage[shortlabels]{enumitem}
\usepackage{amsmath}
\usepackage{amssymb}
\usepackage{amsthm}
\usepackage{mathtools}
\usepackage{xcolor}
\usepackage[margin=1in]{geometry}
\usepackage{hyperref}
\usepackage{cleveref}
\hypersetup{
  colorlinks=true,
  linkcolor=blue,
  filecolor=magenta,
  urlcolor=cyan
}

\newtheorem{theorem}{Theorem}
\newtheorem{thm}{Theorem}[section]
\newtheorem{corollary}[thm]{Corollary}
\newtheorem{lemma}[thm]{Lemma}
\newtheorem{proposition}[thm]{Proposition}
\newtheorem{claim}[thm]{Claim}
\newtheorem{remark}[thm]{Remark}
\newtheorem{question}[thm]{Question}

\newcommand{\whp}{whp}
\newcommand{\defi}{\operatorname{def}}
\newcommand{\odd}{\operatorname{odd}}
\newcommand{\iso}{\operatorname{iso}}
\newcommand{\cN}{\mathcal{N}}
\newcommand{\cO}{\mathcal{O}}
\newcommand{\cE}{\mathcal{E}}
\newcommand{\cK}{\mathcal{K}}
\newcommand{\cP}{\mathcal{P}}
\newcommand{\cQ}{\mathcal{Q}}

\newcommand{\footremember}[2]{%
  \footnote{#2}%
  \newcounter{#1}%
  \setcounter{#1}{\value{footnote}}%
}

\newcommand{\ag}[1]{\textcolor{purple}{(Anna: #1)}}
\newcommand{\sd}[1]{\textcolor{red}{(Sahar: #1)}}
\crefname{originalequation}{original equation}{original equations}
\Crefname{originalequation}{Original equation}{Original equations}
\makeatletter
\newlabel{eq:expansion-log}{{1.1}{source}{}{sourcepointer-1}{}}
\newlabel{eq:expansion-log@cref}{{[originalequation][0][]1.1}{[1][0][]source}{}{}{}}
\makeatother
\crefname{originalequation}{original equation}{original equations}
\Crefname{originalequation}{Original equation}{Original equations}
\makeatletter
\newlabel{eq:mu-range}{{4.2}{source}{}{sourcepointer-2}{}}
\newlabel{eq:mu-range@cref}{{[originalequation][0][]4.2}{[1][0][]source}{}{}{}}
\makeatother
\crefname{originallemma}{original lemma}{original lemmas}
\Crefname{originallemma}{Original lemma}{Original lemmas}
\makeatletter
\newlabel{lem:isolates-far}{{5.1}{source}{}{sourcepointer-3}{}}
\newlabel{lem:isolates-far@cref}{{[originallemma][0][]5.1}{[1][0][]source}{}{}{}}
\makeatother
\crefname{originallemma}{original lemma}{original lemmas}
\Crefname{originallemma}{Original lemma}{Original lemmas}
\makeatletter
\newlabel{lem:component-gap}{{5.2}{source}{}{sourcepointer-4}{}}
\newlabel{lem:component-gap@cref}{{[originallemma][0][]5.2}{[1][0][]source}{}{}{}}
\makeatother
\title{With high probability the random subgraph has one component containing every non-isolated vertex, and its isolated vertices are far apart}
\author{Sahar Diskin, Anna Geisler}
\date{arxiv:2404.14020v4; CC BY 4.0}
\begin{document}
\maketitle
\noindent\textit{Editorial note.} All delivered mathematical text is the author's own, copied from the pinned author source sources/combinatorics-fl04-vspace-1-5cm-perfect-matching-in-p/pinned-original-source.tex; no statement or proof step was rewritten, reordered or removed.\par\medskip
\noindent\textit{Editorial note.} Preamble provenance: the source pins the publisher class scrartcl, which the delivery wrapper does not accept, so this card uses sources/combinatorics-fl04-vspace-1-5cm-perfect-matching-in-p/retained-preamble.tex, which supplies the amsart baseline and then carries the author's own newtheorem and macro declarations (pinned source lines 30-57) verbatim. No macro was invented and no mathematical symbol was changed.\par\medskip
\noindent\textit{Editorial note.} Bibliography provenance: the source carries no thebibliography environment and ends with bibliographystyle abbrv plus bibliography perc (author source lines 1691-1692), so the bibliography is resolved with BibTeX from the identical copy staged with the same arXiv version under sources/product-graph-matchings-2404.14020v4/perc.bib; no entry was edited and the author's own abbrv style is used. The counted proof itself contains no citation.\par\medskip
\noindent\textit{Editorial note.} Selection provenance: this card was opened under the file-based control contract (CONTROL.md), with the mechanical pre-card gate run on its statement and proof interval before the card was written and reporting the interval free.\par\medskip
\noindent\textit{Editorial note.} External references are navigation pointers only; every deliberately kept step of the author's argument stays inside the counted spans.\par\medskip
\noindent\textit{Allowed standard prerequisites (editorial).}\begin{itemize}
\item The two-round exposure of the random subgraph with probabilities p\_1 and p\_2 multiplying to the correct complement, and the range of p displayed as eq:mu-range. Introduced at the start of the counted proof (source lines 542-549); the range is the author's display eq:mu-range at source line 479. Lets the proof apply the already established separation and gap statements to the first round and then merge the large components in the second round.
\item The separation statement for isolated vertices and the gap statement for component orders. The author's lemmas at source lines 489-491 and 514-516, both invoked inside the counted proof at source line 550 and again at 553. Provide the two properties that the proof conditions on before the second round of sprinkling.
\item The expansion estimate used to compare the two rounds. The author's displayed inequality eq:expansion-log at source line 224, invoked inside the counted proof. Shows that the expectation of the first round is the same as the original one up to a factor tending to one.
\end{itemize}\medskip
\ifdefined\hypertarget\hypertarget{sourcepointer-1}{}\fi
\noindent\textit{Source reference.} Original equation 1.1: sources/combinatorics-fl04-vspace-1-5cm-perfect-matching-in-p/pinned-original-source.tex, source lines 224--224. Navigation pointer only: the counted proof invokes this displayed expansion estimate when comparing the two rounds; the display is not reproduced here.\par\medskip
\ifdefined\hypertarget\hypertarget{sourcepointer-2}{}\fi
\noindent\textit{Source reference.} Original equation 4.2: sources/combinatorics-fl04-vspace-1-5cm-perfect-matching-in-p/pinned-original-source.tex, source lines 479--479. Navigation pointer only: the counted proof invokes the author's displayed range for p; the display is not reproduced here.\par\medskip
\ifdefined\hypertarget\hypertarget{sourcepointer-3}{}\fi
\noindent\textit{Source reference.} Original lemma 5.1: sources/combinatorics-fl04-vspace-1-5cm-perfect-matching-in-p/pinned-original-source.tex, source lines 489--489. Navigation pointer only: the counted proof invokes the author's separation statement for isolated vertices; that lemma has its own card in this lane.\par\medskip
\ifdefined\hypertarget\hypertarget{sourcepointer-4}{}\fi
\noindent\textit{Source reference.} Original lemma 5.2: sources/combinatorics-fl04-vspace-1-5cm-perfect-matching-in-p/pinned-original-source.tex, source lines 514--514. Navigation pointer only: the counted proof invokes the author's gap statement for component orders; that lemma has its own card in this lane.\par\medskip
\begin{proposition}\label{prop:unique-component}
Whp, $G_p$ has a unique component containing every non-isolated vertex.
Its isolated vertices are pairwise at distance at least three in $G$.
\end{proposition}

\begin{proof}
Let $p_2=d^{-2}$, and choose $p_1$ so that
\[
  (1-p_1)(1-p_2)=1-p.
\]
The union of independent copies of $G_{p_1}$ and $G_{p_2}$ has the distribution of $G_p$, and
\[
  n(1-p_1)^d=\frac{n(1-p)^d}{(1-d^{-2})^d}=(1+o(1))n (1-p)^d.
\]
The proofs of \Cref{lem:isolates-far,lem:component-gap} hold in the extended regime given in \eqref{eq:mu-range} and thus they apply to $G_{p_1}$.
Also, the expected number of isolated vertices in $G_{p_1}$ is at most $2d^{1/50}$, and so Markov's inequality shows that fewer than $n/2$ vertices are isolated, and hence that there is a non-isolated component, with probability $1-o(1)$.

Using \Cref{lem:isolates-far,lem:component-gap} we may condition on the event that every component of $G_{p_1}$ is either an isolated vertex or has order at least $d^{20}$, and that the isolated vertices are pairwise at distance at least three in $G$. We will show that after sprinkling with $p_2$, typically all the large components, i.e., the components of order at least $d^{20}$ merge.
Let $L$ be the union of the components of order at least $d^{20}$.
Consider a non-trivial partition $A\sqcup B=L$ that respects these components and, by interchanging $A$ and $B$, assume that $a\coloneqq |A|\le|B|$.
Put
\[
  A'=(A\cup N_G(A))\setminus B,
  \qquad B'=V\setminus A'.
\]
Both $A'$ and $B'$ have size at least $a$.
Indeed, they contain $A$ and $B$, respectively.
Let $D$ be the smaller of $A'$ and $B'$.
Then $a\le|D|\le n/2$, and the expansion from \eqref{eq:expansion-log} gives
\[
  e_G(A',B')=e_G(D,D^c)
  \ge \frac{e}{M}|D|\ln\frac{n}{|D|}
  \ge \alpha_Ma,
\]
where $\alpha_M$ is a constant depending solely on $M$.
For the last inequality, use the concavity of $x\mapsto x\ln(n/x)$ on $[a,n/2]$ and evaluate it at the endpoints.

For each edge of $E_G(A',B')$, choose one path of length at most three from $A$ to $B$ containing that edge.
To see this, note that $V\setminus L$ consists of isolated vertices of $G_{p_1}$, no two of which are adjacent in $G$.
An isolated vertex in $A'$ has a neighbour in $A$.
An isolated vertex in $B'$ has no neighbour in $A$; all its $G$-neighbours lie in $L$ because distinct isolated vertices are non-adjacent, and therefore it has a neighbour in $B$.
A fixed edge of $G$ belongs to at most $O(d^2)$ of the chosen paths, because the boundary edge that determines such a path is at distance at most two from it in the line graph of $G$; thus each path has at most three edges.
Greedily retaining such a path and deleting all other paths that share an edge with it therefore leaves at least
\[
  \frac{\alpha_M a}{d^2}
\]
edge-disjoint paths from $A$ to $B$ for some constant $\alpha_M>0$.
The probability that none of them is present when exposing the edges of the second round is at most
\[
  (1-p_2^3)^{\alpha_Ma/d^2}\le \exp\left\{-\frac{\alpha_Ma}{d^8}\right\}.
\]

It remains to take a union bound over the partitions of the large components in $G_{p_1}$. There are at most $N\coloneqq n/d^{20}$ large components.
For a partition whose smaller union has size $a$, that union contains at most $a/d^{20}$ components.
Consequently, the number of possible partitions is at most
\[
  \sum_{j\le a/d^{20}}\binom{N}{j}
  \le \exp\left\{O\left(\frac{a\ln n}{d^{20}}\right)\right\}.
\]
This bound follows from the standard binomial estimate when $a/d^{20}\le N/2$, and from $2^N\le\exp(2a/d^{20})$ otherwise.
Since $\ln n=O_M(d)$, summing over $d^{20}\le a\le n/2$ gives
\[
  \sum_{a=d^{20}}^{n/2}
  \exp\left\{O_M\left(\frac{a}{d^{19}}\right)-\frac{\alpha_Ma}{d^8}\right\}=o(1).
\]
Thus all non-isolated components of $G_{p_1}$ merge after exposing the edges of the second round.

Finally, an isolated vertex of $G_{p_1}$ that receives an edge in the second round cannot join another such vertex, because these vertices are not adjacent in $G$.
It therefore joins the merged component.
The vertices that receive no sprinkled edge remain isolated and by \Cref{lem:isolates-far} they are still of distance at least three.
\end{proof}

\bibliographystyle{plain}
\bibliography{perc}
\end{document}
